\documentclass{article}
\title{MAT1105 ukesoppgaver KMG}
\author{borgebj}
\date{}

\usepackage[norsk]{babel}
\usepackage{mathtools}
\usepackage{amssymb}
\usepackage{graphicx}
\usepackage[a4paper, top=1in, bottom=1in, left=1in, right=1in]{geometry}
\usepackage{subcaption}
\usepackage{float}
\usepackage{enumerate}
\usepackage{dirtytalk}
\usepackage{colortbl}
\usepackage{xcolor} 
\usepackage{amsmath}
\geometry{margin=1in} % Adjust margin as needed



\setlength{\parskip}{2em}
\setlength{\parindent}{0pt}


\begin{document}

\maketitle

\section*{Kapittel 4}
\subsection*{Oppgave 4.4}

La A og B være mengder, og la P være utsanget \(A \subseteq B\)

\( A \subseteq B = \) \\ \say{A er en delmengde av B} eller \say{Alt som er i A er også i B}. \\ \\
Det vil da si at \( \neg P \) = \\ \say{A er ikke en delmengde av B} eller \say{Ikke alt i A også er med i B} \\


Uttrykker følgende utsagn \(\neg P\) ?

\begin{enumerate}[(a)]
    \item \textbf{\(B \subseteq A\)}
    \begin{enumerate}[-]
        \item \textbf{Nei},
        \item \( (B \subseteq A) \quad = \quad (A \subseteq B) \quad \neq \quad \neg (A \subseteq B)\)
    \end{enumerate}
    \item \textbf{\(B \in A\)}
    \begin{enumerate}[-]
        \item \textbf{Nei},
        \item At B er i A er ikke det samme som å si at A ikke er delmengde av B.
    \end{enumerate}
    \item \textbf{\(A = \emptyset\)}
    \begin{enumerate}[-]
        \item \textbf{Nei}, 
        \item Hvis A er det tomme settet, så er P alltid sann.
        \item Dette er fordi den tomme mengden er delmengde av alle mengder.
        \item Derfor holder ikke \(\neg P\)
    \end{enumerate}
    \item \textbf{\(B = \emptyset\)}
    \begin{enumerate}[-]
        \item Svaret avhenger av mengden A.
        \item Hvis \(A = \emptyset\) så er svaret \textbf{Nei}, fordi \(P = \emptyset (A) \subseteq \emptyset(B)\). Og da gjelder ikke \(\neg P\)
        \item Hvis \(A \neq \emptyset\) så er svaret \textbf{Ja}, fordi en ikke-tom delmengde A kan ikke være en delmengde av den tomme mengden B. Og da gjelder \(\neg P\)
    \end{enumerate}
    \newpage
    \item \textbf{\textit{Det finnes en x $\in A \textit{ slik at } x \in B$}}
    \begin{enumerate}[-]
        \item \textbf{Nei},
        \item Samme som å si: \say{Det finnes et element som er med i både A og B}
        \item \(A \subseteq B \text{ og dermed } \neg (A \subseteq B)\) sier noe om hvordan \underline{hele} mengdene A og B oppfører seg.
        \item Denne påstanden beskriver et tilfelle hvor det finnes et element som er i begge, men sier ingenting om hele mengden A.
    \end{enumerate}
    \item \textbf{\textit{Det finnes en x $\in A \textit{ slik at } x \notin B$}}
    \begin{enumerate}[-]
        \item \textbf{Ja},
        \item Samme som å si: \say{Det finnes et element som er med i A, men ikke i B}
        \item Denne påstanden bekrefter at A ikke er en delmengde av B, og dermed at \textbf{\(\neg P er sann\)}
    \end{enumerate}
\end{enumerate}

\subsection*{Oppgave 4.6}
Hvorfor er ikke \( (P \wedge Q) \vee R \text{ og } P \wedge (Q \vee R) \text{ ikke like. }\)

\begin{enumerate}[-]
    \item Hovedsakelig så er de ikke forskjellige på grunn av den ytre konnektiven. 
    \item I utsagnslogisk setninger som disse er det alltid det ytterste konnektivet. 
    \item I første tilfellet så er $\vee$ - \say{eller} konnektivet ytterst. 
\end{enumerate}

\begin{enumerate}[-]
    \item I andre tilfellet så er $\wedge$ - \say{og} konnektivet ytterst.
    \item Vi kan se hvordan dette spiller en rolle i sannhetsverditabellene under, hvor sanne tilordninger er farget grått:

\end{enumerate}

\begin{minipage}[t]{0.48\textwidth}
    \centering
    % Section 1
    \begin{tabular}{c c c | c c c c c}
        $P$ & $Q$ & $R$ & $(P$ & $\wedge$ & $Q)$ & $\vee$ & $R$ \\
        \hline
        0 & 0 & 0 & 0 & 0 & 0 & 0 & 0 \\
        \rowcolor{lightgray}
        0 & 0 & 1 & 0 & 0 & 0 & 1 & 1 \\
        0 & 1 & 0 & 0 & 0 & 1 & 0 & 0 \\
        \rowcolor{lightgray}
        0 & 1 & 1 & 0 & 0 & 1 & 1 & 1 \\
        1 & 0 & 0 & 1 & 0 & 0 & 0 & 0 \\
        \rowcolor{lightgray}
        1 & 0 & 1 & 1 & 0 & 0 & 1 & 1 \\
        \rowcolor{lightgray}
        1 & 1 & 0 & 1 & 1 & 1 & 1 & 0 \\
        \rowcolor{lightgray} 
        1 & 1 & 1 & 1 & 1 & 1 & 1 & 1 \\
    \end{tabular}
    
    \vspace{1em}

    \underline{Sann når:}
    \begin{enumerate}[(i)]
    \centering
        \item P = 0, Q = 0, R = 1
        \item P = 0, Q = 1, R = 1
        \item P = 1, Q = 0, R = 1
        \item P = 1, Q = 1, R = 0
        \item P = 1, Q = 1, R = 1
    \end{enumerate}
\end{minipage}
\hspace{1em} % Horizontal space between the minipages
\begin{minipage}[t]{0.48\textwidth}
    \centering
    
    % Section 2
    \begin{tabular}{c c c | c c c c c}
        $P$ & $Q$ & $R$ & $P$ & $\wedge$ & $(Q$ & $\vee$ & $R)$ \\
        \hline
        0 & 0 & 0 & 0 & 0 & 0 & 0 & 0 \\
        0 & 0 & 1 & 0 & 0 & 0 & 1 & 1 \\
        0 & 1 & 0 & 0 & 0 & 1 & 1 & 0 \\
        0 & 1 & 1 & 0 & 0 & 1 & 1 & 1 \\
        1 & 0 & 0 & 1 & 0 & 0 & 0 & 0 \\
        \rowcolor{lightgray}
        1 & 0 & 1 & 1 & 1 & 0 & 1 & 1 \\
        \rowcolor{lightgray} 
        1 & 1 & 0 & 1 & 1 & 1 & 1 & 0 \\
        \rowcolor{lightgray} 
        1 & 1 & 1 & 1 & 1 & 1 & 1 & 1 \\
    \end{tabular}

    \vspace{1em}

    \underline{Sann når:}
    \begin{enumerate}[(i)]
    \centering
        \item P = 1, Q = 0, R = 1
        \item P = 1, Q = 1, R = 0
        \item P = 1, Q = 1, R = 1
    \end{enumerate}
\end{minipage}

\newpage

\subsection*{Oppgave 4.8}
La \textit{x} være et tall. 
\[\text{\say{Hvorfor er $2x-3 = 5 \Longleftrightarrow x = 4$ Sann?}}\] \\

\( A \Longleftrightarrow B\) betyr at A og B er \textit{ekvivalente}. \\
Det vil si at påstanden kun er sann om enten både A og B er sann, eller om både A og B er usann.
\\ \\
Dette utsagnet er sant, fordi både venstre- og høyresiden er sanne.
\begin{enumerate}[]
    \item Venstresiden: \quad \quad \ \(2x - 3 = 5 \quad \rightarrow \quad 2x = 8 \quad \rightarrow \quad x = 4\)
    \item Høyresiden: \quad \quad \quad \(x = 4\)
\end{enumerate}

I tilfellet hvor enten både venstresiden og høyresiden er usanne, vil hele påstanden også være sann, fordi det er slik ekvialens fungerer. \\ \\
Dersom kun én av sidene var sann, ville hele påstanden vært usann. \\
For eksempel hadde det gitt mening at påstanden var usann hvis vi sa: \(\)

\subsection*{Oppgave 4.11}

\(\neg((\neg P \vee Q) \wedge (\neg Q \vee R))\)

De Morgans lover forteller at:
\begin{enumerate}
    \item $\neg (P \vee Q) \Longleftrightarrow \neg P \wedge \neg Q$
    \item $\neg (P \wedge Q) \Longleftrightarrow \neg P \vee \neg Q$
\end{enumerate}

Ved hjelp av disse kan vi transformere uttrykket videre:

\begin{enumerate}
    \item $\neg ( \neg (P \wedge \neg Q) \wedge \neg (Q \wedge \neg R) )$
    \item $ (P \vee \neg Q) \vee (Q \wedge \neg R)$
    \item $\neg (\neg P \wedge Q) \vee \neg (\neg Q \vee R)$
    \item $\neg ((\neg P \wedge Q) \wedge (\neg Q \vee R))$
\end{enumerate}

Vi kan se at uttrykket kan bli mange og forskjellige.  \\
Tar vi utgangspunkt i for eksempel det andre uttryket, nr 2, så står det: \\ \\
\quad \say{Det er slik at enten (P eller ikke Q), eller (Q og ikke R)} \\ \\

Denne setningen sier ikke mye i seg selv, hovedsakelig fordi utsagnsvariablene våre P, Q og R ikke er tildelt noen utsagn.. \\ \\
Lager vi en sannhetsverditabell for det originale uttrykket kan vi se litt av tilordningene som gjør den sann: 

\hspace{7.4em} 

\(\neg((\neg P \vee Q) \wedge (\neg Q \vee R))\)

\begin{tabular}{c c c | c c c c c c c c c c}
    $P$ & $Q$ & $R$ & $\neg$ & $((\neg$ & $P$ & $\vee$ & $Q)$ & $\wedge$ & $(\neg$ & $Q$ & $\vee$ & $R))$\\
    \hline
    0 & 0 & 0 & 0 & 1 & 0 & 1 & 0 & 1 & 1 & 0 & 1 & 0 \\
    0 & 0 & 1 & 0 & 1 & 0 & 1 & 0 & 1 & 1 & 0 & 1 & 1 \\
    \rowcolor{lightgray}
    0 & 1 & 0 & 1 & 1 & 0 & 1 & 1 & 0 & 0 & 1 & 0 & 0 \\
    0 & 1 & 1 & 0 & 1 & 0 & 1 & 1 & 1 & 0 & 1 & 1 & 1 \\
    \rowcolor{lightgray}
    1 & 0 & 0 & 1 & 0 & 1 & 0 & 0 & 0 & 1 & 0 & 1 & 0 \\
    \rowcolor{lightgray}
    1 & 0 & 1 & 1 & 0 & 1 & 0 & 0 & 0 & 1 & 0 & 1 & 1 \\
    \rowcolor{lightgray}
    1 & 1 & 0 & 1 & 0 & 1 & 1 & 1 & 0 & 0 & 1 & 0 & 0 \\
    1 & 1 & 1 & 0 & 0 & 1 & 1 & 1 & 1 & 0 & 1 & 1 & 1 \\
\end{tabular}

\underline{Sann når:}
\begin{enumerate}[(i)]
    \item P = 0, Q = 1, R = 0
    \item P = 1, Q = 0, R = 0
    \item P = 1, Q = 1, R = 1
\end{enumerate}

\newpage
\subsection*{Oppgave 4.14}

Negasjon av hvert følgende uttrykk, samt tolkninger:

\begin{enumerate}[(a)]
    \item \textbf{\say{Alle sauer har hvit ull}}    
    \begin{enumerate}[-]
        \item \say{Det finnes en som er sau og som ikke har ull.}
        \item La S(x) stå for "x er en sau"
        \item La U(x) stå for "x har ull"
        \item $\exists x (S(x) \wedge \neg I(x))$
    \end{enumerate}
    \item \textbf{\say{Det finnes mennesker som ikke liker is til dessert}}
    \begin{enumerate}[-]
        \item \say{Alle mennesker like is til dessert}
        \item La M(x) stå for "x er et menneske"
        \item La L(x) stå for "x liker is til dessert"
        \item $\forall x (M(x) \wedge L(x))$
    \end{enumerate}
    \item \textbf{\say{Det regner alltid i Bergen}}
    \begin{enumerate}[-]
        \item \say{Det finnes en dag det ikke regner i Bergen}
        \item La R(x) stå for "det regner i Bergen på dag x"
        \item $\exists x \neg Rx$
    \end{enumerate}
    \item \say{Alle myntsamlere har en romersk mynt i samlingen}
    \begin{enumerate}[-]
        \item \say{Det finnes en myntsamler som ikke har en romersk mynt}
        \item La R(y) stå for "y er en romersk mynt"
        \item La M(x) stå for "x er en myntsamler"
        \item La H(x, y) stå for "x har mynten y"
        \item $\exists x (M(x) \land \forall y (R(y) \rightarrow \neg H(x, y)))$
        \item Blir noe lignende: \say{det finnes en x, hvor x er en myntsamler, og for alle y, hvis y er en romersk mynt så har ikke x y}
    \end{enumerate}
    \item \textbf{\say{Noen, men ikke alle, svensker har blondt når}}
    \begin{enumerate}[-]
        \item La S(x) stå for "x er svensk"
        \item La B(x) stå for "x er blond"
        \item Samme som å si $\exists x (S(x) \land B(x))$
        \item Med reglene for negasjon får vi:
        \item \say{Det finnes ingen som er både svensk og blond}
        \item $\forall x \neg(S(x) \land B(x))$
    \end{enumerate}
\end{enumerate}

\newpage

\subsection*{Oppgave 4.15}
\say{Vis at $P \Leftrightarrow Q$ er det samme som $(P \Leftrightarrow Q) \land (\neg P \Rightarrow \neg Q)$}

En måte å vise dette på er å sette opp sannhetsverditabellene for begge: \\
(jeg har fargelagt både de sanne radene grå, og det ytterste konnektivet så det skal være enklere å se)

\begin{minipage}[t]{0.48\textwidth} 
    \begin{tabular}{c c | c >{\columncolor[gray]{0.8}}c c}
    $P$ & $Q$ & $P$ & $\Leftrightarrow$ & $Q$ \\
    \hline
    \rowcolor{lightgray}
     0 & 0 & 0 & 1 & 0 \\
     0 & 1 & 0 & 0 & 1\\
     1 & 0 & 1 & 0 & 0 \\
     \rowcolor{lightgray}
     1 & 1 & 1 & 1 & 1\\
    \end{tabular} 
    \\ 
    \begin{enumerate}
    \item P = 0, Q = 0
    \item P = 1, Q = 1
    \end{enumerate}
\end{minipage}
\hspace{0.5em}
\begin{minipage}[t]{0.48\textwidth} 
    \begin{tabular}{c c | c c c >{\columncolor[gray]{0.8}}c c c c c c }
    $P$ & $Q$ & $(P$ & $\Leftrightarrow$ &  $Q)$ & $\land$ & $(\neg$ & $P$ & $\Rightarrow$ & $\neg$ & $Q)$\\
    \hline
    \rowcolor{lightgray}
     0 & 0 & 0 & 1 & 0 & 1 & 1 & 0 & 1 & 1 & 0  \\
     0 & 1 & 0 & 0 & 1 & 0 & 1 & 0 & 0 & 0 & 1 \\
     1 & 0 & 1 & 0 & 0 & 0 & 0 & 1 & 1 & 1 & 0 \\
     \rowcolor{lightgray}
     1 & 1 & 1 & 1 & 1 & 1 & 0 & 1 & 1 & 0 & 1\\
    \end{tabular} 
    \\ 
    \begin{enumerate}
    \item P = 0, Q = 0
    \item P = 1, Q = 1
    \end{enumerate}
\end{minipage}

Begge uttrykk viser seg å være sanne på samme tilordninger av sannhetsverdier, nemlig når både P og Q er usann, og når både P og Q er sann.


\end{document}
